\documentclass{article}
\usepackage[margin=2cm]{geometry}
\usepackage{amsmath}
\usepackage{graphicx}
\usepackage{booktabs}
\usepackage[utf8]{inputenc}
\usepackage{ragged2e}
\setlength{\emergencystretch}{3em}
\usepackage{paracol}
\begin{document}
\sloppy

\noindent\hfill
\includegraphics[height=1.8em]{figure_header_logo.png}\par\noindent\hrule\vspace{4pt}

\begin{paracol}{2}
\RaggedRight

pared to natural language processing.


\subsection*{11l.NMETHODOLOGY}


\subsection*{A. LDATASETDESCRIPTION}


We use the UCI Human Activity RecognitionUsing Smartphones dataset in its original form1].
This benchmark contains recordings from 30 volunteers aged119-48years who performed1 six
commoncdailyactivities while carryingg a waistmounted Samsung Galaxy S II :smartphone. The
recorded activities include three dynamic motions
(walking, climbing stairs, and descending stairs)
and three static postures (sitting, standing, and
lying).


Each data sample includes 128 consecutive readings collected at 50 Hz, corresponding to a 2.56second window. Nine sensorchannelsare provided: triaxial body aacceleration withgravity removed, triaxial angular velocity fromnthegyroscope, and triaxial total acceleration including
gravity. Together, these channels offer complementary information describing both linear and rotational motion.


The dataset is officially cdivided into 7,352 sampless (70\%) for trainingg and 2,947 samples(30\%)
fortesting. A ssubject-independent split ensures
thatno individual appears mboth sets, which
evaluates modelgeneralization across users rather
than memorization off person-specific：movement
patterns.


\subsection*{B. FPREPROCESSINGPIPELINE}


Raw sensor cdata are loaded directlywithoutt any
additional filtering or external normalization. We
intentionally avoid extra preprocessing to guarantee that observed performancedifferences arise
purely from the modelarchitectures rather: than
preprocessing effects. Batch normalizationn layers
within each network automatically standardize feature distributions during training.


Input tensors are structured as(N,T,C), where
V is the batchsize, T = 128 is the sequence
length, and C = 9 represents the numberof
channels. For convolutional architectures thatTrequire a channel-first layout, inputs are transposed
to (N, C, T) before computation.

% [footer: 4]
\switchcolumn
\RaggedRight

\begin{center}
\includegraphics[width=0.8\linewidth]{figure_001.png}
\end{center}


\textit{FIGURE 1. Ouenview ol four neural architeclures evalualed: (A) CNN·LSTM
baseline.(B) CNN-LSTM with temporal allention,(C) Bidirectional LSTM with
altention.and(D)Translormerencoder.Arrows indicatedalallow direction.}


\subsection*{C. ARCHITECTURE SPECIFICATIONS}


Figure 1 presentsthe four architecturesevaluated
in this study. Lach model embodies a distinct
designphilosophyfor handlingsequential sensor
data.


\subsection*{1)Model A: Convolutional LSTM Baseline}


The first architecture combines convolutional feature extraction with LSTM-based temporal modeling, following the successful DeepConvLSTM
paradigm established by Ord6nez [6]. We deliberately chose two one-dimensional convolutional
blocks to progressively abstract local patterns: the
first block applies 64 filters of size 6 (corresponding to 120ms at 50Hz sampling) followed by ReLU
activation and max pooling with factor 2. This kernel size was selected to capture approximately one
gait cycle during walking activities. The second
block expands to 128 filters of size 3 with similar
activation and pooling.


After convolution, a dropoutt layer with probability (0.1 provides regularization-a conservative
choice given the moderate dataset size. We specifically iutilized a single-layer: ILSTMwith a hidden c dimension of 128, a design derived from the
baseline established by prior work [[24], to process
the compressed features. This hidden size balances

% [footer: VOLUME4.2016]
\end{paracol}

\end{document}
